% part: intuitionistic-logic
% chapter: introduction
% section: syntax

\documentclass[../../../include/open-logic-section]{subfiles}

\begin{document}

\olfileid{int}{int}{syn}

\olsection{अंतःप्रज्ञावादी तर्कशास्त्राची विन्यासमीमांसा}

अंतःप्रज्ञावादी तर्कशास्त्राचा विन्यास विधानीय
तर्कशास्त्राच्या विन्यासासारखाच आहे. अभिजात विधानीय
तर्कशास्त्रात काही संयोजकांची व्याख्या इतर
संयोजकांनी करता येते. उदाहरणार्थ, $!A \lif !B$ ची
$\lnot !A \lor !B$ ने, किंवा $!A \lor !B$ ची
$\lnot(\lnot !A \land \lnot !B)$ ने व्याख्या करता
येते. म्हणून अभिजात तर्कशास्त्राच्या मांडणीत काही
संयोजक अशा व्याख्यांची संक्षिप्त रूपे म्हणून येतात.
अंतःप्रज्ञावादी तर्कशास्त्रात मात्र दोन अपवाद सोडता
तसे होत नाही: $\lnot !A$ हे $!A \lif \lfalse$ चे
संक्षिप्त रूप म्हणून परिभाषित करता येते---आणि
सहसा तसे करतात. अर्थात मग $\lfalse$ ची स्वतःची
व्याख्या इतरांनी केलेली नसली पाहिजे!{} तसेच
$!A \liff !B$ ची अभिजात तर्कशास्त्राप्रमाणे
$(!A \lif !B) \land (!B \lif !A)$ अशी व्याख्या
करता येते.

विधानीय अंतःप्रज्ञावादी तर्कशास्त्रातील !!^{formula}s
ही \emph{!!{propositional variable}s} आणि विधानीय
स्थिरांक~$\lfalse$ यांपासून \emph{तार्किक संयोजके}
वापरून तयार होतात. आपल्याकडे पुढील गोष्टी आहेत:

\begin{enumerate}
\item !!{propositional variable}s चा !!^a{denumerable} संच~$\PVar$: $\Obj p_0$,
  $\Obj p_1$, \dots
  \item !!{falsity} दर्शवणारा विधानीय स्थिरांक~$\lfalse$.
\item तार्किक संयोजके: $\land$
  (संयोग), $\lor$ (विकल्पयोग), $\lif$ (सशर्त)
\item विरामचिन्हे: (, ), आणि स्वल्पविराम.
\end{enumerate}

\begin{defn}[सूत्र]
\ollabel{defn:formulas} विधानीय अंतःप्रज्ञावादी
तर्कशास्त्रातील \emph{!!{formula}s} चा संच~$\Frm[L_0]$
पुढीलप्रमाणे आगमनाने परिभाषित करतो:
\begin{enumerate}
\item $\lfalse$ हे अण्विक !!{formula} आहे.

\item प्रत्येक !!{propositional variable}~$\Obj p_i$ हे अण्विक
  !!{formula} आहे.

\item $!A$ आणि $!B$ ही !!{formula}s असतील, तर $(!A \land
  !B)$ हे !!a{formula} आहे.

\item $!A$ आणि $!B$ ही !!{formula}s असतील, तर $(!A \lor !B)$
  हे !!a{formula} आहे.

\item $!A$ आणि $!B$ ही !!{formula}s असतील, तर $(!A \lif !B)$
  हे !!a{formula} आहे.

\tagitem{limitClause}{याखेरीज दुसरे काहीही !!a{formula} नाही.}{}
\end{enumerate}
\end{defn}

वर दिलेल्या मूलभूत संयोजकांखेरीज पुढील
\emph{परिभाषित} चिन्हेही वापरतो: $\lnot$ (नकार)
आणि $\liff$ (!!{biconditional}). परिभाषित
संकारकांनी तयार झालेल्या सूत्रांचा अर्थ पुढीलप्रमाणे:
\begin{enumerate}
\item $\lnot !A$ हे $!A \lif \lfalse$ चे संक्षिप्त रूप आहे.

\item $!A \liff !B$ हे $(!A \lif !B) \land (!B
  \lif !A)$ चे संक्षिप्त रूप आहे.
\end{enumerate}

$\lnot$ ला औपचारिकरीत्या संक्षिप्त रूप मानले असले,
तरी कधीकधी ते मूलभूत असल्यासारखे~$\lnot$ चे स्पष्ट नियम
आणि व्याख्यांतील खंड देऊ. याचा मुख्य उद्देश सरावाचे
प्रश्न मांडता यावेत हा आहे.

\end{document}
